\par\Needspace{6\baselineskip}\medskip
\hypertarget{native-143-article104}{}\pdfbookmark[2]{104.}{article-104}
\begin{center}104.\end{center}
Superest, ut ex $z$, $n$ vel $c$ elementa ipsa determinemus. Statuendo
\[
\alpha\BareRoot{(ee-1)}=\beta,
\]
\noindent habebitur ex aequatione 6 art. 99
\[
\mathrm{18)}\qquad\beta=\frac{\sin f\mathpunct{.}\BareRoot{rr'}}{\operatorname{tang}2n}\mathpunct{\cdot}
\]
\noindent Combinando hanc formulam cum 12, $12^*$ art. 99, eruitur
\[
\begin{aligned}
\mathrm{19)}\qquad&\BareRoot{(ee-1)}=\operatorname{tang}\psi=\frac{\operatorname{tang}f\operatorname{tang}2n}{2(l-z)}\\[.6em]
\mathrm{19^*)}\qquad&\operatorname{tang}\psi=-\frac{\operatorname{tang}f\operatorname{tang}2n}{2(L+z)},
\end{aligned}
\]
\noindent unde excentricitas commode atque exacte computatur; ex $\beta$ et $\BareRoot{(ee-1)}$ prodibit per divisionem $\alpha$, per multiplicationem $p$, ita ut sit
\[
\begin{aligned}
\alpha&=\frac{2(l-z)\cos f\mathpunct{.}\BareRoot{rr'}}{\operatorname{tang}2n^2}
=\frac{2mm\cos f\mathpunct{.}\BareRoot{rr'}}{yy\operatorname{tang}2n^2}\\[.7em]
&=\frac{kktt}{4yyrr'\cos f^2\operatorname{tang}2n^2}\\[.7em]
&=\frac{-2(L+z)\cos f\mathpunct{.}\BareRoot{rr'}}{\operatorname{tang}2n^2}
=\frac{-2MM\cos f\mathpunct{.}\BareRoot{rr'}}{YY\operatorname{tang}2n^2}\\[.7em]
&=\frac{kktt}{4YYrr'\cos f^2\operatorname{tang}2n^2}
\end{aligned}
\]
\[
\begin{aligned}
p&=\frac{\sin f\mathpunct{.}\operatorname{tang}f\mathpunct{.}\BareRoot{rr'}}{2(l-z)}
=\frac{yy\sin f\mathpunct{.}\operatorname{tang}f\mathpunct{.}\BareRoot{rr'}}{2mm}\\[.7em]
&=\left(\frac{yrr'\sin2f}{kt}\right)^2\\[.7em]
&=\frac{-\sin f\mathpunct{.}\operatorname{tang}f\mathpunct{.}\BareRoot{rr'}}{2(L+z)}
=\frac{-YY\sin f\mathpunct{.}\operatorname{tang}f\mathpunct{.}\BareRoot{rr'}}{2MM}\\[.7em]
&=\left(\frac{Yrr'\sin2f}{kt}\right)^2\mathpunct{\cdot}
\end{aligned}
\]
\noindent\SourceLeaf{144}Expressio tertia et sexta pro $p$, quae omnino identicae sunt cum formulis 18, $18^*$ art. 95, ostendunt, ea quae illic de significatione quantitatum $y$, $Y$ tradita sunt, etiam pro hyperbola valere.
\par
E combinatione aequationum 6, 9 art. 99 deducitur
\[
(r'-r)\BareRoot{\frac{ee-1}{rr'}}=e\sin f\mathpunct{.}\left(C-\frac1C\right);
\]
\noindent introducendo itaque $\psi$ et $\omega$, statuendoque $C=\operatorname{tang}(45^0+N)$, fit
\[
\mathrm{20)}\qquad\operatorname{tang}2N=\frac{2\sin\psi\operatorname{tang}2\omega}{\sin f\cos2\omega}\mathpunct{\cdot}
\]
\noindent Invento hinc $C$, habebuntur valores quantitatis in art. 21 per $u$ expressae pro utroque loco; dein fiet per aequationem III art. 21
\[
\begin{aligned}
\operatorname{tang}\tfrac12v&=\frac{C-c}{(C+c)\operatorname{tang}\tfrac12\psi}\\[.6em]
\operatorname{tang}\tfrac12v'&=\frac{Cc-1}{(Cc+1)\operatorname{tang}\tfrac12\psi}
\end{aligned}
\]
\noindent sive introducendo pro $C$, $c$ angulos $N$, $n$
\[
\begin{aligned}
\mathrm{21)}\qquad&\operatorname{tang}\tfrac12v=\frac{\sin(N-n)}{\cos(N+n)\operatorname{tang}\tfrac12\psi}\\[.6em]
\mathrm{22)}\qquad&\operatorname{tang}\tfrac12v'=\frac{\sin(N+n)}{\cos(N-n)\operatorname{tang}\tfrac12\psi}\mathpunct{\cdot}
\end{aligned}
\]
\noindent Hinc determinabuntur anomaliae verae $v$, $v'$, quarum differentia cum $2f$ comparata simul calculo confirmando inserviet.
\par
Denique per formulam XI art. 22 facile deducitur, intervallum temporis a perihelio usque ad tempus loco primo respondens esse
\[
=\frac{\alpha^{\frac32}}k\left\{\frac{2e\cos(N+n)\sin(N-n)}{\cos2N\cos2n}-\log\operatorname{hyp}\frac{\operatorname{tang}(45^0+N)}{\operatorname{tang}(45^0+n)}\right\}
\]
\noindent et perinde intervallum temporis a perihelio usque ad tempus loco secundo respondens
\[
=\frac{\alpha^{\frac32}}k\left\{\begin{aligned}&\frac{2e\cos(N-n)\sin(N+n)}{\cos2N\cos2n}\\[.5em]&-\log\operatorname{hyp}\operatorname{tang}(45^0+N)\operatorname{tang}(45^0+n)\end{aligned}\right\}.
\]
\noindent Si itaque tempus primum statuitur $=T-\tfrac12t$, adeoque secundum $=T+\tfrac12t$, fit
\[
\mathrm{23)}\qquad T=\frac{\alpha^{\frac32}}k\left\{\frac{e\operatorname{tang}2N}{\cos2n}-\log\operatorname{hyp}\operatorname{tang}(45^0+N)\right\},
\]
\noindent\SourceLeaf{145}unde tempus transitus per perihelium innotescet; denique
\[
\mathrm{24)}\qquad t=\frac{2\alpha^{\frac32}}k\left\{\frac{e\operatorname{tang}2n}{\cos2N}-\log\operatorname{hyp}\operatorname{tang}(45^0+n)\right\},
\]
\noindent quae aequatio, si placet, ad ultimam calculi confirmationem adhiberi potest.
